% Generated by report/tools/tables.py from the result files of the repository. Do not edit.
\begin{table}[tbp]
\centering
\footnotesize
\setlength{\tabcolsep}{3.5pt}
\caption[RL sur grands modèles contre RL classique, premier tirage 2026]{RL appliqué aux grands modèles de langage (famille \texttt{llm} du critère X5) contre RL classique, dans le premier tirage aléatoire de 50 articles de 2026 (ICLR 2026, ICML 2026, NeurIPS 2025) : 27 articles \texttt{llm}, 14 \texttt{classic}, 9 \texttt{theory} (non montrés, comptés dans « Tous »). $k/n$ sur les articles à expériences où le critère s'applique, part en pourcentage entre parenthèses ; intervalle de Wilson à 95~\% pour la colonne « Tous », seul calculé par le script. Seule la grille v1 a ce découpage par famille dans les résultats (\file{study/audit\_stats.json}, harmonisation de \file{study/overrides.json}) : la grille v2 n'est pas ventilée par famille, et aucun test entre familles n'a été préregistré ; le tableau est descriptif. Les deux dernières colonnes rappellent l'échantillon classique élargi à 50 articles (tableau~\ref{tab:sondage-pineau}). Les 18 articles \texttt{llm} et 3 articles \texttt{theory} lus dans la file A3 restent à part : le tri sur le titre en avait déjà écarté beaucoup, ils ne forment pas un échantillon. Sources : \file{study/audit\_stats.json}, \file{study/extension\_stats.json}.}
\label{tab:sondage-lrl}
\begin{tabularx}{\linewidth}{@{}l>{\raggedright\arraybackslash}Xrrrrrr@{}}
\toprule
 & & \multicolumn{4}{c}{Premier tirage 2026 (grille v1)} & \multicolumn{2}{c}{\makecell{RL classique\\{}2026 élargi}} \\
\cmidrule(lr){3-6}\cmidrule(lr){7-8}
Code & Critère & \makecell[r]{RL sur\\{}grands modèles} & \makecell[r]{RL\\{}classique} & Tous & \makecell[r]{\% tous\\{}[Wilson]} & $k/n$ & \% \\
\midrule
Q3 & Tous les hyperparamètres de la méthode & 17/27 (63) & 11/14 (79) & 30/45 & 67 [52\,;\,79] & 36/50 & 72 \\
Q4 & Tous les hyperparamètres des références & 5/27 (19) & 2/14 (14) & 8/44 & 18 [10\,;\,32] & 14/49 & 29 \\
Q5 & Code en lien & 19/27 (70) & 11/14 (79) & 30/45 & 67 [52\,;\,79] & 37/50 & 74 \\
Q6 & Choix des hyperparamètres décrit & 6/27 (22) & 6/14 (43) & 12/45 & 27 [16\,;\,41] & 24/50 & 48 \\
Q7 & Évaluation sur conditions tenues à part & 26/27 (96) & 5/13 (38) & 31/44 & 70 [56\,;\,82] & 22/49 & 45 \\
Q8 & Test statistique appliqué & 2/27 (7) & 0/14 (0) & 2/45 & 4 [1\,;\,15] & 6/50 & 12 \\
X2 & Incertitude rapportée & 4/27 (15) & 14/14 (100) & 21/45 & 47 [33\,;\,61] & 44/50 & 88 \\
X3 & Seuls les meilleurs essais gardés & 1/27 (4) & 1/14 (7) & 2/45 & 4 [1\,;\,15] & 7/50 & 14 \\
X4 & Références réglées avec le même soin & 0/27 (0) & 2/14 (14) & 2/44 & 5 [1\,;\,15] & 8/49 & 16 \\
X7 & Calcul chiffré & 12/27 (44) & 8/14 (57) & 21/45 & 47 [33\,;\,61] & 29/50 & 58 \\
\bottomrule
\end{tabularx}
\end{table}
